\begin{frame}{죄수의 딜레마: MARL을 배우는 2$\times$2 게임}
  경찰이 두 용의자를 \textbf{따로} 심문 --- 각자 \kb{C(침묵,
  coop)}와 \kb{D(자백, defect)} 중 하나. 보상 = 클수록
  좋게 변환:
  \vskip0.5em
  \small
  \begin{tabular}{@{}lcc@{}}
    \toprule
     & 에이전트 2: C (침묵) & 에이전트 2: D (자백) \\
    \midrule
    \textbf{에이전트 1: C (침묵)} & (3, 3) & (0, 5) \\
    \textbf{에이전트 1: D (자백)} & (5, 0) & (1, 1) \\
    \bottomrule
  \end{tabular}
  \vskip0.8em
  {\footnotesize 역사: 1950년경 Merrill Flood·Melvin Dresher가
  RAND에서 실험을 설계하며 만들었고, ``죄수의 딜레마''라는
  이름은 Albert W. Tucker가 붙임. 1928년 John von Neumann가
  2인 제로섬 게임의 ``최적 전략'' 존재를 증명한 것이 게임이론
 의 시작이라면, 죄수의 딜레마는 \textbf{제로섬이 \textbf{아닌}}
  게임이 얼마나 다른지를 보여준 첫 예제.}
  \vskip0.4em
  \begin{itemize}
    \item \textbf{① 지배전략이 D다} --- 상대가 C든 D든 내가
      D가 더 큼(C면 5$>$3, D면 1$>$0): ``상대의 전략을 알아야
      고르는'' 게 아니라 \textbf{``무조건 D''}가 개인 합리적
    \item \textbf{② 그런데 (D,D)는 (C,C)보다 둘 다에게 나쁨}
      --- 둘 다 개인적으로 합리적이면 (1,1)에, 둘 다 C이면
      (3,3)을 나눌 수 있었음
  \end{itemize}
  \vskip0.4em
  {\footnotesize MARL의 핵심 긴장 --- ``\textbf{각자 합리적인
  학습의 결과가 공동으로 최적이지 않을 수 있다}''를 최소
  크기(2$\times$2)로 응축.}
\end{frame}
